%% LyX 2.4.4 created this file.  For more info, see https://www.lyx.org/.
%% Do not edit unless you really know what you are doing.
\documentclass[english]{article}
\usepackage[T1]{fontenc}
\usepackage[latin9]{inputenc}
\usepackage{enumitem}
\usepackage{amsmath}
\usepackage{amsthm}
\usepackage{amssymb}
\usepackage{cancel}
\usepackage{geometry}
\geometry{verbose,tmargin=1in,bmargin=1in,lmargin=1in,rmargin=1in}
\PassOptionsToPackage{normalem}{ulem}
\usepackage{ulem}

\makeatletter
%%%%%%%%%%%%%%%%%%%%%%%%%%%%%% Textclass specific LaTeX commands.
\numberwithin{equation}{section}
\numberwithin{figure}{section}
\newlength{\lyxlabelwidth}      % auxiliary length 

%%%%%%%%%%%%%%%%%%%%%%%%%%%%%% User specified LaTeX commands.
\usepackage{fancyhdr}
\usepackage{lastpage}
\pagestyle{fancy}
\usepackage{enumitem}
\usepackage{ifthen}
\everymath=\expandafter{\the\everymath\displaystyle}
%\DeclareMathSizes{10}{10}{10}{10}
\usepackage{multicol}

\usepackage{tikz}
\usetikzlibrary{patterns}
\usetikzlibrary{plotmarks}
\usepackage{pgfplots}
\pgfplotsset{compat=newest}

%\usepackage{advdate}    % Advancing/saving dates
\usepackage{datetime}   % Dates formatting
%\usdate
%\usepackage{datenumber} % Counters for dates


%\usepackage{calc}
%\usepackage[nomessages]{fp}
% Added by lyx2lyx
\setlength{\parskip}{\smallskipamount}
\setlength{\parindent}{0pt}

\makeatother

\usepackage{babel}
\begin{document}
\renewcommand{\author}{Dr.\,James Ashe}

\newcommand{\classNum}{335}

\newcommand{\classSec}{A}

\newcommand{\examType}{Final Exam}

\renewcommand{\date}{\formatdate{17}{5}{2013}}

%%First page's header%%%%%%%%%%%%%%%%%%%%%%
\fancypagestyle{firststyle} %%header settings for first pagr
{
	\fancyhead[L]{MTH \classNum{}(\classSec{}) \author{}\\
	\textbf{\examType{}} \date{}}
	\fancyhead[C]{\textbf{Name:}}
	\setlength{\headsep}{14pt}
}
%%%%%%%%%%%%%%%%%%%%%%%%%%%%%%%%%%%%%%%%%%

%%Next page's header%%%%%%%%%%%%%%%%%%%%%%%
\setlist{nolistsep}
\lhead{\classNum{}(\classSec{})} 
\chead{\textbf{}} 
\cfoot{\thepage{}~of~\pageref{LastPage}} 
\setlength{\headsep}{5pt} 
%%%%%%%%%%%%%%%%%%%%%%%%%%%%%%%%%%%%%%%%%%%

%%Various settings; see the preamble for other settings%%%%%
%\setenumerate[0]{leftmargin=12pt,itemindent=0pt,labelwidth=0pt}
\setlist{topsep=.5pt}
\renewcommand{\labelenumi}{\textbf{\arabic{enumi}.}}
\renewcommand{\labelenumii}{\textbf{\alph{enumii}.}}
\thispagestyle{firststyle}
%%%%%%%%%%%%%%%%%%%%%%%%%%%%%%%%%%%%%%%%%%

\newcommand{\dspace}{\vspace{.75cm}}
\begin{enumerate}
\item Let $A=\left\{ 1,2,3,5\right\} $ and $B=\left\{ a,b,2,c,5\right\} $

\begin{enumerate}
\item Find $A\cap B$.\dspace
\item Find $A\cup B$.\dspace
\item Find $B-A$.\dspace
\item Find $A-B$.\dspace
\end{enumerate}
\item Let $A$ be a nonempty subset of $B$. Which of the following is true?

\begin{enumerate}
\item [\textbf{(a)}]$A\cap B=B$.
\item [\textbf{(b)}]$A\cup B=A$.
\item [\textbf{(c)}]$A\cup B=B$.
\item [\textbf{(d)}]$A\cap B=\emptyset$.\dspace
\end{enumerate}
\item Let $A=\left\{ 3,5\right\} $ and $B=\left\{ a,b,5\right\} $. 

\begin{enumerate}
\item List all possible subsets of $B$.\dspace\dspace
\item Find $A\times B$. \dspace\dspace
\end{enumerate}
\item Prove that $A\cup\left\{ \emptyset\right\} =A$.\vfill
\item Consider the map $f:\mathbb{R}\times\mathbb{R}\rightarrow\mathbb{R}$
where $f(x,y)=xy$. Prove that $f$ is onto.\vfill\newpage
\item Define the relation $\mathcal{R}$ on the real number set $\mathbb{R}$
as $x\mathcal{R}y\iff xy\geq0$. 

\begin{enumerate}
\item Show that $\mathcal{R}$ is reflexive and symmetric.\dspace\dspace\vfill
\item Is $\mathcal{R}$ an equivalence relation? Explain.\dspace\vfill
\end{enumerate}
\item Let $X$ be a nonempty set. Define the following relation on the set
of subsets of $X$ as follows $A\mathcal{R}B\iff A\cap B\neq\emptyset$
for $A,B\subseteq X$. 

\begin{enumerate}
\item Show that $\mathcal{R}$ is symmetric.\dspace\dspace\vfill
\item Show that $\mathcal{R}$ is not reflexive.\dspace\vfill
\end{enumerate}
\item Let $f:A\rightarrow B$ and $g:B\rightarrow C$ be two functions.

\begin{enumerate}
\item If $g\circ f$ is one-to-one then 

\begin{enumerate}
\item [\textbf{(a)}]$f$ is onto.
\item [\textbf{(b)}]$f$ is one-to-one.
\item [\textbf{(c)}]$g$ is one-to-one.
\item [\textbf{(d)}]Both $f$ and $g$ are one-to-one.
\item [\textbf{(e)}]None of the above. \dspace
\end{enumerate}
\item If $\left|A\right|>\left|B\right|$ and $A\neq\emptyset$ can $f$
be one-to-one?\dspace\newpage
\end{enumerate}
\item Let $S_{3}$ be the symmetric group on three elements. 

\begin{enumerate}
\item Which of the following is true?

\begin{enumerate}
\item [\textbf{(a)}]$S_{3}$ has at least one non-abelian subgroup.
\item [\textbf{(b)}]All proper subgroups of $S_{3}$ are abelian.
\item [\textbf{(c)}]$S_{3}$ is abelian.
\item [\textbf{(d)}]$S_{3}$ is cyclic.\dspace
\end{enumerate}
\item What is the order of $S_{3}$ and $S_{4}$?\dspace\dspace
\end{enumerate}
\item Let $f:G_{1}\rightarrow G_{2}$ be a group homomorphism and let $H\leq G_{2}$.
Recall that $f^{-1}(H)=\left\{ g\in G_{1}|f(g)\in H\right\} $. 

\begin{enumerate}
\item Which of the following is true?

\begin{enumerate}
\item [\textbf{(a)}]$f^{-1}(H)$ consists only of the identity of $G_{1}$.
\item [\textbf{(b)}]$f^{-1}(H)$ \uline{is} a subgroup of $G_{1}$.
\item [\textbf{(c)}]$gf^{-1}(H)=f^{-1}(H)g$ for all $g\in G_{1}$ (the
left and right cosets are equal).
\item [\textbf{(d)}]$f^{-1}(H)$ \uline{is not} a subgroup of $G_{1}$\dspace
\end{enumerate}
\item Let $a,b\in G_{1}$. Does $f(ab)=f(a)f(b)$? \dspace\dspace
\end{enumerate}
\item Let $G$ be a finite group of order $12$ and $g\in G$. 

\begin{enumerate}
\item The order of $g$ \uline{cannot }be:

\begin{enumerate}
\item [\textbf{(a)}]8
\item [\textbf{(b)}]6
\item [\textbf{(c)}]4
\item [\textbf{(d)}]3\dspace
\end{enumerate}
\item Let $\left|H\right|=6$. Does Lagrange's theorem imply that $H\leq G$?\dspace\dspace
\end{enumerate}
\item Let $G$ be a group of order $4$. Which of the following is true?

\begin{enumerate}
\item [\textbf{(a)}]$G$ is a cyclic group.
\item [\textbf{(b)}]$G$ is $\mathbb{Z}_{2}\times\mathbb{Z}_{2}$.
\item [\textbf{(c)}]$G$ is $\mathbb{Z}_{4}$.
\item [\textbf{(d)}]$G$ is abelian.\dspace\newpage
\end{enumerate}
\item Let $G=\left\{ e,a_{1},a_{2},a_{3},a_{4},a_{5}\right\} $ be a group
where $e$ is the identity.

\begin{enumerate}
\item Find a subgroup of $G$.\dspace\dspace
\item Find a subgroup of $G$ distinct from the group given for part \textbf{(a)}.\dspace\dspace
\item Can $\left|a_{2}\right|=5$?\dspace
\end{enumerate}
\item Let $H$ and $K$ be subgroups be non-trivial subgroups of $G$. Prove
or disprove that $H\cap K$ is a subgroup of $G$.\vfill
\item Let $R$ be a ring and $a,b\in R$. 

\begin{enumerate}
\item Which statement is true?

\begin{enumerate}
\item [\textbf{(a)}]If $ab=0$ then $a=0$ or $b=0$.
\item [\textbf{(b)}]If $a^{2}=0$ then $a=0$.
\item [\textbf{(c)}]If $a\neq0$ and $b\neq0$ then $ab\neq0$.
\item [\textbf{(d)}]None of the above.\dspace
\end{enumerate}
\item Show that $\mathbb{Q}$, the set of rational numbers, is a ring under
addition and multiplication. \vfill
\end{enumerate}
\end{enumerate}

\end{document}
